% Generated draft: requires mathematical and source review.
\documentclass[11pt]{article}
\usepackage[margin=25mm]{geometry}
\usepackage{fontspec}
\setmainfont{Latin Modern Roman}
\usepackage{amsmath,amssymb,mathtools,unicode-math}
\setmathfont{Latin Modern Math}
\usepackage{enumitem,needspace,xurl,hyperref}
\setlength{\parindent}{0pt}
\setlength{\parskip}{5pt}
\setlength{\emergencystretch}{3em}
\newcommand{\N}{\mathbb{N}}
\newcommand{\Z}{\mathbb{Z}}
\newcommand{\Q}{\mathbb{Q}}
\newcommand{\R}{\mathbb{R}}
\newcommand{\C}{\mathbb{C}}
\newcommand{\F}{\mathbb{F}}
\newcommand{\A}{\mathbb{A}}
\newcommand{\PP}{\mathbb{P}}
\newcommand{\E}{\mathbb{E}}
\newcommand{\eps}{\varepsilon}
\newcommand{\dd}{\,\mathrm{d}}
\begin{document}
\section*{UnsolvedMath ID: 30001513}
\textbf{Source problem code:} catalogue entry
\subsection{Definitions and mathematical statement}
Let \(n\in\Z\) with \(n\ge 2\), and let \(F_n\) be the free group of rank \(n\). Let \(\phi:F_n\to F_n\) be an injective endomorphism that is proper, understood here to mean \(\phi(F_n)\subsetneq F_n\), and irreducible in the sense intended by the supplied source. The record does not define irreducibility, so this hypothesis remains unresolved; in particular, no choice between conventions concerning invariant free factors, free-factor systems, or iterates is imposed here.

Define the mapping-torus group, equivalently the ascending HNN extension associated to \(\phi\), by
\[
G_\phi=\left\langle F_n,t\ \middle|\ t^{-1}gt=\phi(g)\text{ for every }g\in F_n\right\rangle.
\]
Also define
\[
\operatorname{SL}_2(\C)=\{A\in M_2(\C):\det A=1\}.
\]
An embedding means an injective group homomorphism; no discreteness condition on its image is specified.

The supplied question, read universally, asks: for every \(n\ge2\) and every proper injective endomorphism \(\phi:F_n\to F_n\) satisfying the source's intended irreducibility hypothesis, does there exist an injective homomorphism
\[
\rho:G_\phi\longrightarrow\operatorname{SL}_2(\C)?
\]
A fully self-contained formulation requires clarification of irreducibility. The interpretation of proper as nonsurjective should also be checked against the cited report.
\subsection{Short English statement}
Does every mapping-torus group of an irreducible, injective, nonsurjective endomorphism of a finite-rank free group of rank at least two admit a faithful representation in \(\operatorname{SL}_2(\C)\)? The source does not define irreducibility, so the exact class of endomorphisms remains to be clarified.
\subsection{Sources}
\begin{itemize}
\item Supplied local UnsolvedMath registry record, permanent ID \(30001513\), title ``Linear Representations of Ascending Free-Group Mapping Tori,'' source problem code \texttt{catalogue entry}, in the Oberwolfach Reports Open Problems collection. No individual proposer is identified. The record gives proposed year 2010.
\item Supplied source citation: \emph{Geometric Group Theory, Hyperbolic Dynamics and Symplectic Geometry} (2011), Oberwolfach Reports, pp.\ 1761--1825, DOI \texttt{10.4171/owr/2010/30}. \url{https://doi.org/10.4171/owr/2010/30}. The supplied background describes these reports as documenting mathematical workshops and their open-problem discussions.
\item Supplied literature source labeled ``Calegari Dunfield'': \url{https://arxiv.org/abs/math/0412136}.
\item Supplied literature source labeled ``original OWR report'': \url{https://doi.org/10.4171/owr/2010/30}.
\item Literature assessment dated exactly \texttt{2026-08-21}: the supplied curation labels the problem open and states that the embedding question remains unresolved for irreducible proper injective free-group endomorphisms, while Calegari and Dunfield only show that irreducibility cannot be omitted. This is supplied dated curation derived from the OWR record, not an independently verified literature assessment; its time-sensitive status and cited claims require independent checking.
\item The registry labels the statement \texttt{corrected\_verified} and reports that its mathematical content was checked against the official Oberwolfach report, with notation, context, and question form normalized without changing the source-level claim. This is a supplied verification assertion, not verification performed for this draft.
\end{itemize}
\end{document}
